\documentclass[aps,prd,twocolumn,showpacs,superscriptaddress,preprintnumbers]{revtex4-2}
\usepackage{amsmath,amssymb,amsfonts}
\usepackage{graphicx}
\usepackage{hyperref}
\usepackage{dsfont}
\usepackage{bm}
\usepackage{booktabs}

\begin{document}

\title{Two Rules, One Universe: Gravity and Quantum Decoherence from Information Capacity}

\author{Huang Zhongchang}
\affiliation{Independent Researcher, Nanjing, China}

\date{\today}

\begin{abstract}
We propose that both gravitational attraction and quantum decoherence are consequences of a single mechanism: the bounded information capacity of physical cells, formalized through only two axioms---causality (A1) and capacity bound (A2)---plus one structural choice for the dynamical variable.
The free-capacity fraction $q$ of each cell drives an information current $J$, whose relaxation yields the \textbf{unified telegraph equation}
$\partial_t^2 q + \gamma_0(1-q)\partial_t q = c^2\nabla^2(\ln q)$,
derived from the fundamental current equation
$\partial_t J_{ij} = -(c^2/a^2)(\ln q_i - \ln q_j) - \gamma_0(1-\bar{q})J_{ij}$
with $\bar{q}=(q_i+q_j)/2$.
For small perturbations $\delta q$ around $q=1$, damping vanishes at linear order, yielding the wave equation $\Box(\delta q)=0$; the static limit gives $\nabla^2(\ln q)=0$ directly, without patching.
The damping $\gamma(q)=\gamma_0(1-q)$ is the simplest ansatz consistent with A2: occupied cells block incoming information flow.
The cell spacing $a$ is the single new dimensional parameter; identifying $a$ with $\ell_P\equiv\sqrt{\hbar G/c^3}$ follows from the naturalness requirement that the theory contain no new scales beyond $c$, $\hbar$, $G$.
In vacuum ($q\to1$), $\gamma\to0$ and the dynamics select the Lorentz-invariant sector $\Box(\delta q)=0$---preferred-frame effects are confined to matter-filled regions, and Lorentz symmetry, while always a subsector of the theory, is \textbf{activated} by the vacuum ground state rather than imposed as an axiom.
The static nonlinear equation $\nabla^2(\ln 1/q)=0$ yields $q(r)=e^{-GM/rc^2}$, recovering Newtonian gravity at large $r$.
We prove an H-theorem for the full nonlinear theory via a Lyapunov functional, establishing the thermodynamic consistency of the framework.
The theory makes two structural predictions: a wave--diffusion crossover in gravitational dynamics at $k=\gamma/(2c)$, and a soft boundary replacing the classical event horizon.
Limitations are assessed honestly, including the tension with the holographic bound and the need for extension to full general relativity.
\end{abstract}

\maketitle

%%%%%%%%%%%%%%%%%%%%%%%%%%%%%%%%%%%%%%%%%%%%%%%%%%%%%%%%%%%%%%%%%%%%%%
\section{Introduction}
\label{sec:intro}
%%%%%%%%%%%%%%%%%%%%%%%%%%%%%%%%%%%%%%%%%%%%%%%%%%%%%%%%%%%%%%%%%%%%%%

Unifying quantum mechanics with gravity. Few problems in fundamental physics have resisted solution for so long.
Quantum field theory---built by quantizing classical fields on a fixed spacetime background---is the most precisely tested framework in physics, yet general relativity, describing gravity as spacetime curvature, has passed every observational test~\cite{Will:2014}.
The two frameworks are structurally incompatible: GR is a classical, deterministic, geometric theory; quantum mechanics is a probabilistic theory on Hilbert space.
String theory~\cite{Green:1987sp,Polchinski:1998rq}, loop quantum gravity~\cite{Rovelli:2004tv,Ashtekar:2004eh}, causal dynamical triangulations~\cite{Ambjorn:2001cv}, asymptotic safety~\cite{Weinberg:1979,Reuter:2012}, entropic gravity~\cite{Verlinde:2010hp,Jacobson:1995ab}, and the quantum memory matrix~\cite{Neukart:2024qmm}---numerous approaches exist.
None has yielded a complete, experimentally verified theory.

This paper takes a different approach. Rather than quantizing geometry or geometrizing quantum fields, we ask a simpler question: can gravity be derived from \emph{information-theoretic} first principles?
Put differently, if both quantum mechanics and gravity are consequences of the same underlying constraint---the finite information capacity of spacetime---then their apparent incompatibility may dissolve.
That is the hypothesis we explore here.

Our starting point is Zilly's (2026) theorem~\cite{Zilly:2026}. From finite information capacity and self-referential consistency (``graded equality''), the full mathematical structure of quantum mechanics---complex amplitudes, Born rule, unitary dynamics---emerges as the unique consistent description.
This is a foundational result. It means quantum mechanics is itself a consequence of bounded capacity, not a separate axiom.
Zilly's theorem also establishes $\hbar$ as the natural scale of the capacity bound, fixing phase space volume quantization~\cite{Zilly:2026}.
We accept this theorem as the quantum-mechanical foundation.
We build upward from it toward gravity. This is not a derivation of quantum mechanics---it is a construction that takes QM as the already-derived starting point and asks what else falls out of the same scaffolding.

We introduce exactly \textbf{two axioms}:

\begin{enumerate}
    \item[A1.] \textbf{Causality exists.} There is asymmetric, irreversible influence between discrete information cells---a directed partial order on events.
    \item[A2.] \textbf{Capacity bounded.} Each cell stores at most 1 bit of information. The maximum information propagation speed is $c$. The state space of $N$ cells is $\mathcal{X}=\{0,1\}^N$.
\end{enumerate}

From these axioms we construct a dynamical theory of the free-capacity fraction $q_i(t)$, the fraction of cell $i$'s 1-bit capacity that remains available for new information.

\textbf{Why 1 bit?} Axiom~A1 guarantees that at least two states of a cell are distinguishable. The minimal distinguishable difference---one bit---is the irreducible unit of information.
A2 imposes that this unit is also the \emph{maximum}: each cell can distinguish at most two states ($s_i\in\{0,1\}$), hence capacity is exactly 1~bit.
The value is not arbitrary; it follows from the conjunction of distinguishability (A1, which demands at least 1~bit) and boundedness (A2, which limits to at most 1~bit).
Higher capacities ($n>1$~bits per cell) would introduce an unexplained integer $n$ and violate the principle that the theory contain no gratuitous structure.

\textbf{Zilly independence.} Throughout this paper, Zilly's (2026) theorem~\cite{Zilly:2026} provides the quantum-mechanical context---assuring us that standard QM is compatible with, and indeed derivable from, finite-capacity principles.
However, \emph{no result in Sections~\ref{sec:dynamics}--\ref{sec:selforg} depends on Zilly's theorem.}
The derivation of the unified telegraph equation, the dynamical Lorentz selection, the static Newtonian limit, and the H-theorem require only A1, A2, and the Tier~2 structural choices enumerated in Table~\ref{tab:tiers}.
Zilly's theorem enters solely to connect the classical gravitational limit back to the quantum-mechanical foundation, closing the conceptual loop from quantum to classical---but the classical derivation stands independently.

The main results are:

\begin{itemize}
    \item The unified current equation with $\bar{q}=(q_i+q_j)/2$ and the continuity equation combine to yield $\partial_t^2 q + \gamma_0(1-q)\partial_t q = c^2\nabla^2(\ln q)$. For small perturbations around $q=1$, $\ln q\approx q-1$ recovers the linear telegraph equation; the $\gamma\to0$ vacuum limit yields $\Box(\ln q)=0$; and the static limit directly gives $\nabla^2(\ln q)=0$.
    \item In vacuum ($q\to1$), $\gamma\to0$ and Lorentz invariance is dynamically selected as the ground state of the empty universe, not imposed as an axiom.
    \item The preferred-frame effects from $\gamma>0$ are confined to matter-filled regions ($q<1$), where they are suppressed by $(1-q)\sim GM/rc^2\sim10^{-6}$ near the solar surface---consistent with all Lorentz-violation constraints.
    \item In the static limit, $\nabla^2(\ln 1/q)=0$, whose spherically symmetric solution $q(r)=e^{-GM/rc^2}$ recovers the Newtonian potential $\Phi=-GM/r$ at large $r$.
    \item The decomposition $\nabla^2 q/q = |\nabla q|^2/q^2$ unifies quantum decoherence (diffusion-dominant at $q\approx1$) and gravitational collapse (self-organization-dominant at $q\approx0$) as two regimes of a single information-dynamic process.
    \item An H-theorem for the full nonlinear theory is proved via a Lyapunov functional, establishing that the dynamics are thermodynamically consistent.
\end{itemize}

This work builds on and differs from prior approaches in specific ways.
Jacobson~\cite{Jacobson:1995ab} derived the Einstein equations from thermodynamic considerations applied to local Rindler horizons; we go further, deriving the Newtonian potential from information dynamics without assuming any metric structure.
Verlinde~\cite{Verlinde:2010hp} proposed entropic gravity; we provide the microscopic mechanism---the $q$-field and $J$-current---that his approach lacks.
Padmanabhan~\cite{Padmanabhan:2010} explored the thermodynamic structure of gravitational field equations; our framework reproduces similar structure from purely information-theoretic axioms and proves an explicit H-theorem.
Bassani and Magueijo~\cite{Bassani:2025htma} showed that diffeomorphism invariance can emerge from a Markov chain of stochastic choices; our Lorentz selection is deterministic and requires far fewer assumptions, though the emergent symmetry is weaker (Lorentz vs.\ diffeomorphism).
We provide a detailed, methodologically fair comparison in Sec.~\ref{sec:discussion}, applying the same tiered-assumption methodology to both theories.

The paper is organized as follows.
Section~\ref{sec:axioms} formalizes the axioms and constructs the cell graph, including a tiered classification of all assumptions.
Section~\ref{sec:dynamics} derives the unified $q$-field dynamics from the fundamental current equation.
Section~\ref{sec:gamma} derives $\gamma(q)=\gamma_0(1-q)$ from the capacity bound and frames $\gamma_0$ through naturalness.
Section~\ref{sec:lorentz} proves that Lorentz invariance is dynamically selected by vacuum.
Section~\ref{sec:static} takes the static limit and recovers Newtonian gravity through the unified derivation.
Section~\ref{sec:selforg} analyzes self-organization, proves the H-theorem, and gives the stability analysis.
Section~\ref{sec:predictions} presents testable predictions.
Section~\ref{sec:discussion} discusses limitations, compares with HTMAU using tiered methodology, and outlines open directions.

%%%%%%%%%%%%%%%%%%%%%%%%%%%%%%%%%%%%%%%%%%%%%%%%%%%%%%%%%%%%%%%%%%%%%%
\section{Axioms and Cell Graph}
\label{sec:axioms}
%%%%%%%%%%%%%%%%%%%%%%%%%%%%%%%%%%%%%%%%%%%%%%%%%%%%%%%%%%%%%%%%%%%%%%

\subsection{Foundational Result}

We take as given Zilly's (2026) theorem~\cite{Zilly:2026}:

\begin{quote}
\textbf{Zilly Theorem (2026).} Let a physical system have finite information capacity $C<\infty$ and satisfy self-referential consistency (``graded equality'' of distinguishable states). Then: (1) the state representation necessarily involves complex numbers; (2) the probability rule is the Born rule $P(a)=|\langle a|\psi\rangle|^2$; (3) the dynamics are unitary; and (4) the resulting theory is standard quantum mechanics~\cite{Zilly:2026}.
\end{quote}

We do not re-derive or question this result.
We accept it as the quantum-mechanical foundation upon which we build the classical gravitational limit.

What we should emphasize: Zilly's theorem is a recent result (2026).
Like any new theoretical development, it awaits independent verification and community scrutiny.
But the architecture of our argument does not crumble if the theorem requires revision---the classical gravitational limit derived in Sections~\ref{sec:static}--\ref{sec:selforg} depends only on A1--A2 and the Tier~2 structural choices, not on the specific mechanism by which quantum mechanics emerges.
The key insight we exploit is that Zilly's finite-capacity argument applies at \emph{every} scale---capacity is bounded per cell, not just globally.
Crucially, Zilly's theorem introduces $\hbar$ as the scale of the capacity bound: the quantization of phase space in units of $\hbar$ follows from the same finite-capacity argument, so $\hbar$ enters our framework not as a separate axiom but as part of the quantum-mechanical foundation provided by Zilly's theorem.

\subsection{Information Cells and Graph}

Consider a set of $N$ discrete information cells arranged on a graph $G=(V,E)$.
Each vertex $i\in V$ is a cell.
Edges $(i,j)\in E$ are potential information transfer channels.
The graph has the following properties:

\begin{enumerate}
    \item \textbf{Geometric embedding.} Cells are embedded in $d=3$ spatial dimensions. We take $d=3$ as an empirical input; the theory does not predict spatial dimensionality.
    \item \textbf{Spacing.} Adjacent cells are separated by a fundamental length scale $a$. We identify $a$ with the Planck length $\ell_P\equiv\sqrt{\hbar G/c^3}$. This identification follows from the naturalness requirement that the theory contain no new scales beyond those already present in $c$, $\hbar$, and $G$ (see Sec.~\ref{sec:gamma} for detailed justification).
    \item \textbf{Locality.} Direct information transfer occurs only between adjacent cells (nearest neighbors on the graph).
    \item \textbf{Speed limit.} Information propagates at most at speed $c$ between adjacent cells, so the minimum time step is $\Delta t = a/c$.
\end{enumerate}

\subsection{Tiered Classification of Assumptions}
\label{sec:tiered}

All inputs to the theory are classified into four tiers, separating irreducible physical principles from modeling choices.

\begin{table}[h]
\centering
\caption{Tiered classification of all assumptions entering the $q$-field framework. Tiers 0--1 are irreducible inputs; Tier~2 choices are protected by universality; Tier~3 quantities are fixed by Tiers~0--2.}
\label{tab:tiers}
\begin{tabular}{p{0.12\columnwidth}p{0.40\columnwidth}p{0.40\columnwidth}}
\toprule
\textbf{Tier} & \textbf{Content} & \textbf{Status} \\
\midrule
Tier 0 & A1: Causality exists (directed partial order). & Irreducible physical axiom. \\
(Physical Axioms) & A2: Capacity bound (1 bit per cell, max speed $c$). & Irreducible physical axiom. \\
\midrule
Tier 1 & $c$: maximum signal speed. & Empirically measured. \\
(Empirical Inputs) & $\hbar$: quantum of action (via Zilly's theorem). & Empirically measured. \\
& $G$: Newton's constant. & Empirically measured. \\
& $d=3$: spatial dimensionality. & Empirically observed. \\
\midrule
Tier 2 & Regular 3D lattice. & Simplest choice; universality protects \\
(Modeling Choices) & Nearest-neighbor edges. & macroscopic predictions. \\
& $\phi(q)=-\ln q \in \Phi$. & Any $\phi\in\Phi$ gives same qualitative \\
& & physics (Sec.~\ref{sec:static}). \\
\midrule
Tier 3 & $\gamma_0 = c/a = c/\ell_P$. & Fixed by Tiers 0--2 + naturalness. \\
(Derived Quantities) & $D_0 = c^2/\gamma_0 = c\ell_P$. & Derived from $\gamma_0$. \\
\bottomrule
\end{tabular}
\end{table}

\textbf{Tier 0--1 are irreducible inputs.} They are the physical facts that any theory must accommodate.
\textbf{Tier 2 choices are protected by universality.} Any potential $\phi\in\Phi=\{\phi(q)\mid\phi'(q)<0,\;\phi(q)\to\infty\;\text{as}\;q\to0\}$ gives the same qualitative physics: the static equation is $\nabla^2\phi(q)=0$, the $1/r$ potential emerges, and the resistance to information inflow diverges as $q\to0$.
The specific choice $\phi(q)=-\ln q$ is the simplest member of this universality class.
\textbf{Tier 3 quantities are fixed by Tiers 0--2.} The damping scale $\gamma_0=c/a$ is the only dimensional parameter constructed from the cell spacing $a$ and the speed limit $c$.
Identifying $a$ with $\ell_P$ is a naturalness argument (Sec.~\ref{sec:gamma}): the theory introduces $a$ as the single new scale, and the only natural candidate consistent with the requirement that no new scales appear beyond $c$, $\hbar$, $G$ is the Planck length.

\subsection{The Free-Capacity Fraction $q$}

Each cell $i$ has a total information capacity of 1 bit (Axiom~A2).
At any time $t$, some fraction of this capacity is occupied (storing information) and some is free (available for new information).
Define:

\begin{equation}
q_i(t) \equiv \frac{\text{free capacity of cell } i}{\text{total capacity of cell } i} \in [0, 1].
\label{eq:qdef}
\end{equation}

The free-capacity fraction $q_i$ is the fundamental dynamical variable of our theory.
It is a real scalar field on the discrete graph.
We will show that in the continuum limit, $q$ obeys equations that reproduce both quantum wave dynamics and classical gravitational field equations, depending on the regime.

\subsection{Physical Interpretation}

\begin{itemize}
    \item $q_i = 1$: cell $i$ is completely empty. All capacity is available. This corresponds to asymptotically flat spacetime far from any mass.
    \item $q_i = 0$: cell $i$ is completely full. No capacity remains. This corresponds to the limiting case of infinite information density---the analog of a gravitational singularity, though we will see that $q=0$ is approached asymptotically, never reached at finite $r$.
    \item $0 < q_i < 1$: intermediate occupation. Gradients in $q$ drive information currents, which manifest as gravitational effects.
\end{itemize}

A crucial point: $q$ is \emph{not} a field in spacetime.
Rather, $q$ is a property of the information substrate from which spacetime itself emerges.
The metric structure we perceive as gravity is a coarse-grained description of $q$-field gradients.

\subsection{Coarse-Graining and the Continuum Limit}

Here is a subtlety.
Axiom~A2 states that each individual cell stores exactly 1~bit ($s_i\in\{0,1\}$), so a single cell's free-capacity fraction can only be 0 or 1.
The field $q_i(t)\in[0,1]$ used throughout this paper is defined on \emph{mesoscopic blocks} containing $N_b\gg 1$ Planck-scale cells, over which the binary occupancy is averaged:
$q_i \equiv N_b^{-1}\sum_{k\in\text{block}(i)}(1-s_k)$.
For $N_b$ sufficiently large, $q_i$ is effectively continuous, justifying the use of PDEs.
The discrete Bernoulli noise from individual binary cells is suppressed by $1/\sqrt{N_b}$ and is negligible at macroscopic scales.

With $N\sim10^{184}$ Planck cells in the observable universe and $N_b\sim10^{60}$ cells per mesoscopic block (e.g., at the nuclear scale), continuum error estimates are $\mathcal{O}(10^{-123})$ (see SM~B5).
All dynamical equations derived below describe this coarse-grained field; the binary nature of A2 manifests only at the Planck scale and does not affect the continuum phenomenology.

%%%%%%%%%%%%%%%%%%%%%%%%%%%%%%%%%%%%%%%%%%%%%%%%%%%%%%%%%%%%%%%%%%%%%%
\section{$q$-Field Dynamics: The Unified Equation}
\label{sec:dynamics}
%%%%%%%%%%%%%%%%%%%%%%%%%%%%%%%%%%%%%%%%%%%%%%%%%%%%%%%%%%%%%%%%%%%%%%

\subsection{Continuity Equation}

Information is conserved in transfer. It moves between cells but is neither created nor destroyed.
For cell $i$, the rate of change of occupied capacity equals the net current into the cell:

\begin{equation}
\frac{\partial}{\partial t}(1 - q_i) = -\sum_{j \in \partial i} J_{ij},
\label{eq:continuity_raw}
\end{equation}

where $\partial i$ denotes the set of cells adjacent to $i$, and $J_{ij}$ is the information current from cell $i$ to cell $j$.
By convention, $J_{ij}>0$ means information flows from $i$ to $j$.
Since $(1-q_i)$ is the occupied fraction, Eq.~\eqref{eq:continuity_raw} simplifies to:

\begin{equation}
\boxed{\frac{\partial q_i}{\partial t} = \sum_{j \in \partial i} J_{ij}}.
\label{eq:continuity}
\end{equation}

The current is antisymmetric: $J_{ij}=-J_{ji}$, reflecting that information leaving $i$ arrives at $j$.
This is the discrete analog of the conservation law $\partial_t q + \nabla\cdot\mathbf{J}=0$ in the continuum.

\subsection{Fundamental Current Equation}

We need a constitutive relation for $J_{ij}$.
The current is driven by gradients in free capacity: information tends to flow from cells with more free capacity to cells with less, as the latter are ``information-rich'' and resist further inflow by the capacity bound.
The simplest local, causal dynamics for the current is a Cattaneo-type relaxation equation~\cite{Cattaneo:1948,Jou:2010}.

\textbf{The nonlinear current equation.}
Three physical constraints must be satisfied:
\begin{enumerate}
    \item The current must vanish when $q_i=q_j$ (no gradient, no flow).
    \item The driving force must diverge as $q_i\to0$ or $q_j\to0$ (no cell can accept information beyond its 1-bit capacity).
    \item The current direction must be downhill in $q$.
\end{enumerate}

These are satisfied by the logarithmic gradient $(\ln q_i-\ln q_j)$, which belongs to the universality class $\Phi$.
For $q_i,q_j\approx1$, $\ln q_i-\ln q_j\approx q_i-q_j$, recovering the linear approximation---but as $q\to0$, the logarithmic gradient diverges, encoding the capacity bound's resistance to further inflow.

Now the damping.
Occupied cells block information flow, so the effective damping seen by a current $J_{ij}$ should scale with the average occupation of the two cells involved.
We therefore write the damping with $\bar{q}\equiv(q_i+q_j)/2$:

\begin{equation}
\boxed{\frac{\partial J_{ij}}{\partial t} = -\frac{c^2}{a^2}(\ln q_i - \ln q_j) - \gamma_0(1-\bar{q})J_{ij}}.
\label{eq:current_fundamental}
\end{equation}

This is the \textbf{fundamental current equation} of the theory.
One equation, used consistently throughout.
The coupling constant $c^2/a^2$ is the natural frequency scale (inverse light-crossing time squared between adjacent cells).
The damping coefficient $\gamma(q)=\gamma_0(1-q)$ will be derived from the capacity bound in Sec.~\ref{sec:gamma}; for now, $\gamma_0$ is a dimensionful constant with units of inverse time.

Equation~\eqref{eq:current_fundamental} satisfies all three constraints: it vanishes at equal $q$, diverges at $q\to0$, and drives flow downhill.
It is \textbf{not} two separate equations (linear and nonlinear) patched together---it is one unified equation valid in all regimes.

\subsection{Unified Telegraph Equation}

Differentiate the continuity condition with respect to time and substitute the current equation.
The algebra is straightforward but the result is compact:

\begin{align}
\frac{\partial^2 q_i}{\partial t^2} &= \sum_{j \in \partial i} \frac{\partial J_{ij}}{\partial t} \nonumber \\
&= \sum_{j \in \partial i} \left[-\frac{c^2}{a^2}(\ln q_i - \ln q_j) - \gamma_0(1-\bar{q}_{ij})J_{ij}\right] \nonumber \\
&= -\frac{c^2}{a^2} \sum_{j \in \partial i} (\ln q_i - \ln q_j) - \gamma_0 \sum_{j \in \partial i}(1-\bar{q}_{ij})J_{ij}.
\label{eq:telegraph_step}
\end{align}

In the continuum limit $a\to0$, the discrete Laplacian of $\ln q$ becomes:
\begin{equation}
\frac{1}{a^2}\sum_{j \in \partial i}(\ln q_i - \ln q_j) \;\longrightarrow\; -\nabla^2(\ln q).
\end{equation}

For the damping sum, when $q$ varies smoothly across the lattice, $\bar{q}_{ij}\approx q_i$ for all neighbors $j$.
Then:
\begin{equation}
\sum_{j \in \partial i}(1-\bar{q}_{ij})J_{ij} \approx (1-q_i)\sum_{j \in \partial i}J_{ij} = (1-q_i)\frac{\partial q_i}{\partial t}.
\end{equation}

We thus obtain the \textbf{unified telegraph equation}:

\begin{equation}
\boxed{\frac{\partial^2 q}{\partial t^2} + \gamma_0(1-q)\frac{\partial q}{\partial t} = c^2\nabla^2(\ln q)}.
\label{eq:telegraph_unified}
\end{equation}

This is the fundamental dynamical equation of $q$-field theory.
In the derivation, we have neglected the cross-term $\sum_j J_{ij}\cdot\nabla(\ln q_i)$, which is $\mathcal{O}(|\nabla q|^2)$ and subdominant to the $(1-q)\partial_t q$ term for small perturbations around uniformity; the full derivation retaining this term is given in SM~A3.3.

\subsection{Limiting Cases}

The unified telegraph equation~\eqref{eq:telegraph_unified} interpolates between qualitatively different regimes, all contained within the single equation.
We examine three of them.

\subsubsection{Small-Amplitude Limit: Linear Telegraph}

For perturbations $\delta q$ around a background $q_0$:
\begin{equation}
\ln q = \ln(q_0 + \delta q) \approx \ln q_0 + \frac{\delta q}{q_0}.
\end{equation}

For the important case of vacuum perturbations ($q_0=1$):
\begin{equation}
\ln(1+\delta q) \approx \delta q \quad\Longrightarrow\quad \nabla^2(\ln q) \approx \nabla^2(\delta q).
\end{equation}

The damping term becomes $\gamma_0(1-1-\delta q)\partial_t(\delta q) = -\gamma_0\delta q\,\partial_t\delta q$, which is \textbf{second order} in $\delta q$.
Therefore, to linear order in perturbations around $q=1$:

\begin{equation}
\boxed{\frac{\partial^2 \delta q}{\partial t^2} - c^2\nabla^2\delta q = 0 \quad\Longrightarrow\quad \Box(\delta q) = 0}.
\label{eq:kg_vacuum}
\end{equation}

Perturbations around vacuum propagate as massless waves at speed $c$.
The damping term, being second order in the perturbation, does not appear at linear order.

More generally, for perturbations around $q_0<1$:
\begin{equation}
\frac{\partial^2 \delta q}{\partial t^2} + \gamma_0(1-q_0)\frac{\partial \delta q}{\partial t} - c^2\nabla^2\left(\frac{\delta q}{q_0}\right) = 0,
\label{eq:linear_telegraph_general}
\end{equation}
which is the linear telegraph equation with effective damping $\gamma_{\rm eff}=\gamma_0(1-q_0)$ and an effective diffusion constant $D_{\rm eff}=c^2/[\gamma_0(1-q_0)]$.

\subsubsection{Vacuum Limit: $\gamma\to0$}

In exact vacuum ($q=1$ everywhere), the damping term vanishes identically:
\begin{equation}
\gamma(1) = \gamma_0(1-1) = 0,
\end{equation}
and Eq.~\eqref{eq:telegraph_unified} reduces to:
\begin{equation}
\boxed{\Box(\ln q) = 0},
\label{eq:box_lnq}
\end{equation}
where $\Box\equiv\partial_t^2-c^2\nabla^2$ is the d'Alembertian.
This is the nonlinear wave equation for $\ln q$ in vacuum.
For small perturbations, it reduces to $\Box(\delta q)=0$ as shown above.

\subsubsection{Overdamped Regime: $\gamma\gg\partial_t$}

The damping function $\gamma(q)=\gamma_0(1-q)$ is bounded above by $\gamma_0$ (attained at $q=0$).
There is \textbf{no} ``$\gamma\to\infty$'' limit---$\gamma$ is always finite, with maximum value $\gamma_0=c/\ell_P\approx1.85\times10^{43}$\,s$^{-1}$.

A physically important regime occurs when the inertial term is negligible compared to damping:
\begin{equation}
|\partial_t^2 q| \ll \gamma|\partial_t q|.
\end{equation}
This condition is met on timescales $t\gg\gamma^{-1}$ or, equivalently, for wavenumbers $k\ll k_{\rm cross}\equiv\gamma/(2c)$.
In this \textbf{overdamped regime}, Eq.~\eqref{eq:telegraph_unified} reduces to the effective diffusion equation:

\begin{equation}
\boxed{\frac{\partial q}{\partial t} = D_{\rm eff}\nabla^2(\ln q), \qquad D_{\rm eff} \equiv \frac{c^2}{\gamma_0(1-q)}}.
\label{eq:diffusion}
\end{equation}

The crossover between wave and diffusive behavior is controlled by the dimensionless parameter:
\begin{equation}
R \equiv \frac{\omega}{\gamma} \quad\text{or}\quad \frac{k}{k_{\rm cross}}.
\end{equation}
$R\gg1$: wave regime (inertial term dominates). $R\ll1$: diffusion regime (damping dominates).
This crossover is a sharp prediction of the theory (Sec.~\ref{sec:predictions}).

\textbf{Rigorous analysis.}
The transition from hyperbolic (telegraph) to parabolic (diffusion) behavior is a singular perturbation problem.
As $\epsilon\equiv\gamma^{-1}\to0$, the highest time derivative is multiplied by a small parameter, creating a boundary layer at $t\sim\epsilon$.
For $t\gg\epsilon$, the outer solution obeys the diffusion equation~\eqref{eq:diffusion}; the inner solution matches the initial wavefront.
This structure is standard for telegraph $\to$ diffusion limits~\cite{Jou:2010} and guarantees that information propagates causally (front speed $c$) even in the diffusion-dominated regime.

%%%%%%%%%%%%%%%%%%%%%%%%%%%%%%%%%%%%%%%%%%%%%%%%%%%%%%%%%%%%%%%%%%%%%%
\section{Damping from Capacity: $\gamma(q)=\gamma_0(1-q)$}
\label{sec:gamma}
%%%%%%%%%%%%%%%%%%%%%%%%%%%%%%%%%%%%%%%%%%%%%%%%%%%%%%%%%%%%%%%%%%%%%%

\subsection{Physical Argument}

In Eq.~\eqref{eq:current_fundamental}, the damping term $-\gamma_0(1-\bar{q})J_{ij}$ represents the loss of coherence in the information current.
The physical origin is straightforward: information propagation requires \emph{free capacity} in the receiving cell.
When a cell is occupied (probability $1-q$), it cannot accept new information and therefore blocks the current.

Within any local neighborhood, the fraction of occupied cells encountered by a current is $(1-q)$.
It follows that:

\begin{equation}
\gamma \propto (1-q).
\label{eq:gamma_scaling}
\end{equation}

Damping is not an external friction parameter.
It is an \emph{intrinsic} consequence of the capacity bound (A2).
The more occupied the cells, the harder it is for information to flow coherently.
That observation is the central physics of this section.

\subsection{The Linear Ansatz}

The simplest functional form consistent with the physical constraints is:

\begin{equation}
\boxed{\gamma(q) = \gamma_0 (1-q)}.
\label{eq:gamma_linear}
\end{equation}

This satisfies:

\begin{itemize}
    \item \textbf{Vacuum limit:} As $q\to 1$ (completely empty cells), $\gamma\to 0$. In vacuum, there are no occupied cells to block information flow, so damping vanishes entirely. This is a dynamical consequence of $q\to 1$, not a limit we take by hand.
    \item \textbf{Saturation limit:} As $q\to 0$ (completely full cells), $\gamma\to\gamma_0$. All cells are occupied, blocking is maximal, and the current experiences the strongest possible damping.
\end{itemize}

\textbf{Honest statement:} The linear form $\gamma(q)=\gamma_0(1-q)$ is a \emph{modeling ansatz}, not a rigorous derivation from A2.
Nonlinear corrections (e.g., $\gamma\propto(1-q)^\alpha$ with $\alpha\neq 1$) may exist.
The essential point---that $\gamma$ must vanish as $q\to 1$---follows from A2 regardless of the specific functional form.
The linear ansatz captures this essential physics with minimal assumptions.

\subsection{Naturalness Argument for $\gamma_0$}
\label{sec:naturalness}

$\gamma_0$ has dimensions of inverse time.
The cell spacing $a$ is the \textbf{single new dimensional parameter} introduced by the theory.
Combined with the speed limit $c$, it defines the natural frequency scale:

\begin{equation}
\gamma_0 = \frac{c}{a}.
\label{eq:gamma0_form}
\end{equation}

What fixes $a$?
The theory contains exactly three dimensionful constants from Tier~1: $c$, $\hbar$, and $G$.
If the theory is to contain \textbf{no new scales} beyond those already present in known physics, then $a$ must be constructed from $c$, $\hbar$, and $G$.
Dimensional analysis gives a unique combination with dimensions of length:

\begin{equation}
a = \ell_P \equiv \sqrt{\frac{\hbar G}{c^3}} \approx 1.62 \times 10^{-35}\;\text{m}.
\label{eq:lp_def}
\end{equation}

Consequently:

\begin{equation}
\boxed{\gamma_0 \equiv \frac{c}{\ell_P} = \sqrt{\frac{c^5}{\hbar G}} \approx 1.85 \times 10^{43}\;\text{s}^{-1}}.
\label{eq:gamma0}
\end{equation}

We stress: this is a \textbf{naturalness argument}, not a logical necessity.
The theory introduces $a$ as a free parameter; we \emph{choose} to identify it with $\ell_P$ because doing so introduces no new scales.
This choice has falsifiable consequences: if observations constrain $\gamma_0 \neq c/\ell_P$, the identification is excluded.

In the language of Tier~3 (Table~\ref{tab:tiers}), $\gamma_0=c/\ell_P$ is a derived quantity, fixed by Tier~0--2 inputs plus the naturalness requirement.

\textbf{Dimensionless prefactor.}
Dimensional analysis alone permits an arbitrary $\mathcal{O}(1)$ dimensionless prefactor: $\gamma_0 = \alpha\,c/a$ with $\alpha\sim\mathcal{O}(1)$.
Setting $\alpha=1$ is the simplest choice and is adopted throughout this paper.
All qualitative predictions (Lorentz selection in vacuum, wave--diffusion crossover existence, Newtonian limit recovery, H-theorem) are independent of $\alpha$.
Naturalness prefers $\alpha\sim\mathcal{O}(1)$; the exact value must be determined by observation, making $\alpha$ the single dimensionless free parameter of the theory.

\subsection{Modified Telegraph Equation}

Substituting $\gamma(q)=\gamma_0(1-q)$ into Eq.~\eqref{eq:telegraph_unified}:

\begin{equation}
\boxed{\frac{\partial^2 q}{\partial t^2} + \gamma_0(1-q)\frac{\partial q}{\partial t} = c^2\nabla^2(\ln q)}.
\label{eq:telegraph_gammaq}
\end{equation}

Key properties:
\begin{enumerate}
    \item \textbf{Nonlinearity:} The damping term $\gamma_0(1-q)\partial_t q$ couples the field to its own time derivative. Additionally, the right-hand side $\nabla^2(\ln q)$ provides nonlinear spatial coupling.
    \item \textbf{Vacuum linearization:} For perturbations $\delta q$ around $q=1$, the damping term is $O(\delta q^2)$ and the spatial term linearizes to $\nabla^2\delta q$, giving $\Box(\delta q)=0$ exactly at linear order.
    \item \textbf{Static nonlinearity:} In the static limit $\partial_t\to0$, Eq.~\eqref{eq:telegraph_gammaq} directly yields $\nabla^2(\ln q)=0$, the equation whose solution gives Newtonian gravity (Sec.~\ref{sec:static}).
\end{enumerate}

%%%%%%%%%%%%%%%%%%%%%%%%%%%%%%%%%%%%%%%%%%%%%%%%%%%%%%%%%%%%%%%%%%%%%%
\section{Lorentz Invariance Dynamically Selected by Vacuum}
\label{sec:lorentz}
%%%%%%%%%%%%%%%%%%%%%%%%%%%%%%%%%%%%%%%%%%%%%%%%%%%%%%%%%%%%%%%%%%%%%%

\subsection{The Selection Theorem}

\begin{quote}
\textbf{Dynamic Lorentz Selection.}
In vacuum, defined as the state $q=1$ (all cells empty), the capacity-dependent damping vanishes: $\gamma(1)=\gamma_0(1-1)=0$.
The dynamics reduce to $\Box(\ln q)=0$ and, for perturbations, $\Box(\delta q)=0$.
Since $\Box$ is invariant under Lorentz transformations, vacuum spacetime is automatically Lorentz invariant.
Lorentz invariance is \textbf{not} imposed as an axiom---it is \textbf{dynamically selected} as the ground state of the empty universe.
\end{quote}

We did not assume Lorentz invariance.
Not as an axiom, not as a Lagrangian symmetry.
The vacuum state of the theory \emph{automatically} possesses Lorentz invariance because the damping---the only term that could break it---vanishes identically at $q=1$.

The mechanism is transparent:
\begin{enumerate}
    \item A2 (capacity bound) implies that information flow is blocked by occupied cells.
    \item This blocking produces damping $\gamma\propto(1-q)$.
    \item In vacuum ($q=1$), there are no occupied cells $\Rightarrow$ no blocking $\Rightarrow$ $\gamma=0$.
    \item With $\gamma=0$, the dynamics are $\Box(\ln q)=0$, which is Lorentz invariant.
\end{enumerate}

\textbf{Honest statement:} Lorentz selection here is weaker than the diffeomorphism emergence claimed in some other frameworks.
Our theory always \emph{contained} a Lorentz-invariant sector ($\Box(\ln q)=0$); the result is that vacuum \emph{automatically selects} this sector rather than us discovering it by taking a limit by hand.
The nontrivial claim is that the preferred-frame sector ($\gamma>0$) is dynamically confined to matter-filled regions.

\subsection{Preferred-Frame Effects Are Confined to Matter}

For $q<1$ (anywhere matter is present), $\gamma>0$ and the dynamics are not Lorentz invariant.
The damping term picks out a preferred rest frame---the frame of the information substrate.
This is physically expected: decoherence requires a preferred frame~\cite{Bassi:2013mta}, and here that frame is the rest frame of the cell graph.

Is the resulting Lorentz violation compatible with existing constraints?
For a spherically symmetric mass $M$ at distance $r$, Eq.~\eqref{eq:q_solution} gives:

\begin{equation}
1 - q(r) = 1 - e^{-GM/rc^2} \approx \frac{GM}{rc^2}.
\label{eq:lv_estimate}
\end{equation}

At the solar surface ($M=M_\odot$, $r=R_\odot$):

\begin{equation}
1 - q \approx \frac{GM_\odot}{R_\odot c^2} \approx 2 \times 10^{-6}.
\label{eq:solar_lv}
\end{equation}

So $\gamma/\gamma_0 \approx 2\times 10^{-6}$ at the solar surface.
All existing Lorentz-violation constraints~\cite{Kostelecky:2008,Mewes:2013} are satisfied.
The theory was not tuned to satisfy them---the same $GM/rc^2$ factor that makes gravity weak also makes Lorentz violation small. In the laboratory ($M\sim$ kg, $r\sim$ m), $1-q\sim 10^{-27}$, utterly negligible.

%%%%%%%%%%%%%%%%%%%%%%%%%%%%%%%%%%%%%%%%%%%%%%%%%%%%%%%%%%%%%%%%%%%%%%
\section{Static Limit and Gravitational Potential}
\label{sec:static}
%%%%%%%%%%%%%%%%%%%%%%%%%%%%%%%%%%%%%%%%%%%%%%%%%%%%%%%%%%%%%%%%%%%%%%

\subsection{Unified Static Derivation}

Consider stationary configurations: $\partial_t q=0$ and $\partial_t J_{ij}=0$.
From Eq.~\eqref{eq:current_fundamental} with $\partial_t J_{ij}=0$:

\begin{equation}
0 = -\frac{c^2}{a^2}(\ln q_i - \ln q_j) - \gamma_0(1-\bar{q}_{ij})J_{ij}.
\label{eq:static_J_eq}
\end{equation}

Solving for the equilibrium current:

\begin{equation}
J_{ij} = -\frac{c^2}{\gamma_0(1-\bar{q}_{ij})a^2}(\ln q_i - \ln q_j).
\label{eq:J_static}
\end{equation}

The zero-net-current condition from Eq.~\eqref{eq:continuity} with $\partial_t q_i=0$ gives $\sum_{j\in\partial i}J_{ij}=0$ for all $i$:

\begin{equation}
\sum_{j\in\partial i} \frac{c^2}{\gamma_0(1-\bar{q}_{ij})a^2}(\ln q_i - \ln q_j) = 0.
\label{eq:zero_net_current}
\end{equation}

For a smoothly varying $q$-field, $\bar{q}_{ij}\approx q_i$ is nearly constant across all neighbors $j$, and the factor $\gamma_0(1-\bar{q}_{ij})$---while finite---cancels from the sum.
The condition reduces to:

\begin{equation}
\sum_{j \in \partial i} (\ln q_i - \ln q_j) = 0.
\label{eq:laplace_ln_disc}
\end{equation}

In the continuum limit $a\to0$:

\begin{equation}
\boxed{\nabla^2(\ln 1/q) = 0}.
\label{eq:laplace_ln}
\end{equation}

This is the fundamental static equation of $q$-field theory, derived directly from the unified dynamics without any patching, separate ``nonlinear correction,'' or ad hoc modification.
The derivation is one continuous argument from Eq.~\eqref{eq:current_fundamental} to Eq.~\eqref{eq:laplace_ln}:

\begin{center}
\small
$\partial_t J_{ij}=0 \;\Rightarrow\; J_{ij} \propto (\ln q_i-\ln q_j)/(1-\bar{q}_{ij}) \;\Rightarrow\; \sum_j J_{ij}=0 \;\Rightarrow\; \sum_j(\ln q_i-\ln q_j)=0 \;\Rightarrow\; \nabla^2(\ln q)=0.$
\end{center}

\textbf{Universality.}
The specific choice $\phi(q)=-\ln q$ belongs to the universality class $\Phi=\{\phi(q) \mid \phi'(q)<0,\;\phi(q)\to\infty\;\text{as}\;q\to0\}$.
Any $\phi\in\Phi$ gives the same qualitative physics: resistance diverges as $q\to0$, the static equation is $\nabla^2\phi(q)=0$, and the $1/r$ potential emerges.
The choice $-\ln q$ is the simplest member of this class and is protected by universality---the macroscopic physics does not depend on microscopic details of $\phi$.

\subsection{Spherically Symmetric Solution}

For spherical symmetry, $\ln(1/q)$ depends only on $r$.
Setting $\nabla^2(\ln 1/q)=0$:

\begin{equation}
\frac{1}{r^2}\frac{d}{dr}\left(r^2 \frac{d}{dr}\ln\frac{1}{q}\right) = 0,
\label{eq:ode}
\end{equation}

which integrates to $r^2\frac{d}{dr}\ln(1/q)=\text{const}$.
Denoting the integration constant by $GM/c^2$ (chosen to match the Newtonian limit):

\begin{equation}
r^2 \frac{d}{dr}\ln\frac{1}{q} = \frac{GM}{c^2}.
\label{eq:first_int}
\end{equation}

Integrating again with the boundary condition $q\to1$ as $r\to\infty$:

\begin{equation}
\boxed{q(r) = \exp\left(-\frac{GM}{rc^2}\right)}.
\label{eq:q_solution}
\end{equation}

This is the central result: the free-capacity fraction around a spherically symmetric mass $M$.

\subsection{Recovery of Newtonian Gravity}

For $r\gg GM/c^2$ (the weak-field regime):

\begin{equation}
q(r) = 1 - \frac{GM}{rc^2} + \mathcal{O}\left(\frac{G^2 M^2}{r^2 c^4}\right).
\label{eq:q_expand}
\end{equation}

Define the gravitational potential $\Phi(r)\equiv-c^2(1-q(r))$.
In the weak-field limit:

\begin{equation}
\boxed{\Phi(r) = -\frac{GM}{r} + \mathcal{O}(r^{-2})}.
\label{eq:newton}
\end{equation}

Newton's law is recovered in the weak-field, static limit.
We now show that the same information-theoretic principles fix the full spacetime metric, not just the Newtonian potential.

\subsection{Emergent Spacetime Metric}

\textbf{Causal locking and $g_{00}$.}
The causal locking hypothesis (Sec.~\ref{sec:lorentz}) states that the local proper time $\tau$ advances at a rate proportional to the free-capacity fraction: $d\tau/dt = q$.
In the rest frame of a static observer, $d\tau^2 = -g_{00}dt^2$, giving $g_{00} = -q^2$.

\textbf{Information propagation and $g_{rr}$.}
The spatial metric follows from how information traverses the cell graph.
Each cell-to-cell hop occurs at speed $c = a/t_P$, but succeeds only if the target cell is free---which occurs with probability $q$.
The mean hop time is $t_P/q$, so the effective propagation speed is $c_{\rm eff} = c q$.
The proper distance traversed by a signal is $dl_{\rm proper} = dl_{\rm coord}/q$, yielding $g_{rr} = 1/q^2$ in isotropic coordinates.

\textbf{Full static metric.}
With $q(r) = e^{-GM/rc^2}$, the emergent line element is:
\begin{equation}
\boxed{ds^2 = -e^{-2GM/rc^2}c^2dt^2 + e^{2GM/rc^2}(dr^2 + r^2d\Omega^2)}.
\label{eq:metric}
\end{equation}

\textbf{PPN parameters and consistency with GR.}
Expanding to first post-Newtonian order ($U \equiv GM/rc^2 \ll 1$):
$g_{00} = -(1 - 2U + 2U^2 + \ldots)$, $g_{rr} = 1 + 2U + 2U^2 + \ldots$.
In the standard PPN gauge, this yields $\beta = 1$ and $\gamma = 1$---identical to GR at leading order.
The perihelion precession, light deflection, and Shapiro time delay all match GR predictions exactly at this order.
Second-order deviations ($\Delta g_{rr} \sim U^2/2$) are $\sim10^{-16}$ for Mercury, eight orders of magnitude below current MESSENGER precision, but may be accessible to future strong-field observations.

\textbf{Honest statement:} We recover the Newtonian limit of GR, \emph{not} the full Einstein equations.
The $1/r$ potential was not put in by hand---it follows from spherical symmetry, the static equation $\nabla^2(\ln 1/q)=0$, and the boundary condition $q\to1$ at infinity.
The constant $GM/c^2$ is an integration constant; we identify $G$ as Newton's constant and $M$ as the source mass.
The $1/r$ dependence and exponential form are derived, not assumed.

\subsection{The Role of $d=3$: Marginal Dimensionality}
\label{sec:d3}

Spatial dimensionality $d=3$ is taken as an empirical input (Tier~1).
But it plays a distinguished role.
In $d$ dimensions, the static equation $\nabla^2(\ln 1/q)=0$ integrates to $\ln(1/q)\propto r^{-(d-2)}$ (for $d\neq2$) or $\ln(1/q)\propto\ln r$ (for $d=2$).
The gravitational potential is then $\Phi\sim r^{-(d-2)}$.

Only $d=3$ yields the $1/r$ force law that permits stable orbits, long-range structure formation, and the scaling $\Phi\propto M/r$ that we observe.
In this sense, $d=3$ is a \textbf{marginal} dimension for long-range forces: $d<3$ gives forces that fall off too slowly for stable bound systems; $d>3$ gives forces that fall off too rapidly for hierarchical structure formation.
The fact that our universe selects $d=3$ is an empirical fact, but the theory highlights why this particular dimension enables the rich gravitational phenomenology we observe.

\textbf{Matter coupling.}
The static equation $\nabla^2(\ln 1/q)=0$ describes vacuum configurations; the spherically symmetric solution $q(r)=e^{-GM/rc^2}$ is obtained with an integration constant identified as $GM/c^2$ to match the Newtonian limit.
A complete theory must derive how matter (occupied information capacity) sources the $q$-field in general configurations, yielding a Poisson-like equation with a source term $\propto(1-q)/a^3$.
The derivation of the correct dimensionful prefactor and the full coupling to the stress-energy tensor is left to future work; the vacuum solution calibrated to the Newtonian limit suffices for the static, weak-field results derived here.
A natural candidate source term, consistent with the identification of occupied capacity as information charge, is $\nabla^2(\ln 1/q) = (4\pi G/c^2)\,\rho_m$ with $\rho_m \propto (1-q)/a^3$, which reduces to the standard Poisson equation in the weak-field limit.
Establishing this coupling from the microscopic current dynamics is a priority for extending the framework beyond the vacuum Newtonian regime.

%%%%%%%%%%%%%%%%%%%%%%%%%%%%%%%%%%%%%%%%%%%%%%%%%%%%%%%%%%%%%%%%%%%%%%
\section{Self-Organization: Unifying Decoherence and Collapse}
\label{sec:selforg}
%%%%%%%%%%%%%%%%%%%%%%%%%%%%%%%%%%%%%%%%%%%%%%%%%%%%%%%%%%%%%%%%%%%%%%

\subsection{Decomposition of the Static Equation}

Expand the static nonlinear equation~\eqref{eq:laplace_ln}:

\begin{equation}
\nabla^2(\ln q) = \frac{\nabla^2 q}{q} - \frac{|\nabla q|^2}{q^2} = 0.
\label{eq:expand_ln}
\end{equation}

This gives:

\begin{equation}
\boxed{\frac{\nabla^2 q}{q} = \frac{|\nabla q|^2}{q^2}}.
\label{eq:decomp}
\end{equation}

Equation~\eqref{eq:decomp} expresses a balance between two competing processes.
On one side, the diffusion term $\nabla^2 q/q$---the standard Laplacian diffusion that smooths gradients and homogenizes $q$.
This drives the system toward uniform $q$, corresponding to decoherence (loss of structured information).
On the other, the self-organization term $|\nabla q|^2/q^2$---always non-negative, amplifying existing gradients.
This is positive feedback driving structure formation.

Both terms come from the \textbf{same equation}.
Not two mechanisms. Two faces of one information-dynamic process.

\subsection{Regime Crossover}

Which term dominates depends on $q$:

\begin{itemize}
    \item \textbf{$q\approx1$ (low information density):} The denominator $q^2$ in the self-organization term is $\approx1$, but $\nabla q$ is small (near-flat capacity). Diffusion dominates, and the system tends toward uniformity. This is the \textbf{quantum regime}: decoherence drives the system toward classicality.

    \item \textbf{$q\approx0$ (high information density):} The self-organization term diverges as $q^{-2}$, overwhelming diffusion. The system amplifies any existing gradient, driving further collapse. This is \textbf{gravitational collapse}: structure formation as runaway information concentration.
\end{itemize}

In the intermediate regime ($0<q<1$), the balance determines stability.
From quantum decoherence at $q\approx1$ through astrophysical accretion to gravitational collapse at $q\to0$---one unified picture of structure formation across all scales.

\subsection{H-Theorem: Lyapunov Functional for the Nonlinear Theory}

\label{sec:lyapunov}

We prove that the full nonlinear dynamics possess an H-theorem, establishing thermodynamic consistency.
Define the generalized energy functional:

\begin{equation}
E[q,\mathbf{J}] = \int d^3x\,\left[\Phi(q) + \frac{1}{2\kappa}|\mathbf{J}|^2\right],
\label{eq:lyapunov_E}
\end{equation}

where $\kappa\equiv c^2/a^2$ and $\Phi(q)$ satisfies $\Phi'(q)=\ln q$. (The continuum field $\mathbf{J}$ used here absorbs a factor of $a$ relative to the discrete edge currents; see SM for the equivalent formulation in terms of the bare discrete currents.)
Explicitly, $\Phi(q)=q\ln q - q + 1$ (the constant is chosen so that $\Phi(1)=0$, the vacuum energy).

The continuum limit of the fundamental current equation~\eqref{eq:current_fundamental} is:

\begin{equation}
\frac{\partial \mathbf{J}}{\partial t} = -\kappa\nabla(\ln q) - \gamma_0(1-q)\mathbf{J}.
\label{eq:J_continuum}
\end{equation}
(The continuum field $\mathbf{J}$ in this section is oriented opposite to the discrete edge currents $J_{ij}$ of Sec.~\ref{sec:dynamics}, a convention chosen to simplify the Lyapunov functional. The H-theorem is independent of this orientation choice.)

Together with the continuity equation $\partial_t q = -\nabla\cdot\mathbf{J}$, we compute the time derivative of $E$:

\begin{align}
\frac{dE}{dt} &= \int d^3x\,\left[\Phi'(q)\frac{\partial q}{\partial t} + \frac{1}{\kappa}\mathbf{J}\cdot\frac{\partial\mathbf{J}}{\partial t}\right] \nonumber \\
&= \int d^3x\,\Bigl[\ln q\,(-\nabla\cdot\mathbf{J}) \nonumber \\
&\qquad + \frac{1}{\kappa}\mathbf{J}\cdot\bigl(-\kappa\nabla(\ln q) - \gamma_0(1-q)\mathbf{J}\bigr)\Bigr] \nonumber \\
&= \int d^3x\,\Bigl[-\ln q\,\nabla\cdot\mathbf{J} - \mathbf{J}\cdot\nabla(\ln q) - \frac{\gamma_0(1-q)}{\kappa}|\mathbf{J}|^2\Bigr].
\label{eq:dEdt_step1}
\end{align}

Integrating the first two terms by parts:

\begin{equation}
\int d^3x\,\bigl[-\ln q\,\nabla\cdot\mathbf{J} - \mathbf{J}\cdot\nabla(\ln q)\bigr] = -\oint_{\partial V} (\ln q)\,\mathbf{J}\cdot d\mathbf{S}.
\label{eq:boundary_term}
\end{equation}

For isolated systems ($q\to1$, $\mathbf{J}\to0$ at spatial infinity) or periodic boundary conditions, the surface term vanishes.
Then:

\begin{equation}
\boxed{\frac{dE}{dt} = -\int d^3x\,\frac{\gamma_0(1-q)}{\kappa}\,|\mathbf{J}|^2 \;\leq\; 0}.
\label{eq:dEdt_final}
\end{equation}

Since $\gamma_0(1-q)\geq0$, $\kappa>0$, and $|\mathbf{J}|^2\geq0$, we have $dE/dt\leq0$ for all configurations.
Equality holds only when $\mathbf{J}=0$ everywhere (static equilibrium) or $q=1$ everywhere (vacuum).

\textbf{This proves the H-theorem for the full nonlinear theory.}
The dynamics are irreversible and thermodynamically consistent.
The Lyapunov functional identifies the static configurations ($\mathbf{J}=0$, $\nabla^2(\ln q)=0$) as the attractors of the dynamics, with the exponential solution~\eqref{eq:q_solution} as the minimum free-energy configuration for a given total information charge (mass).

The physical picture is clean.
$\int\Phi(q)\,d^3x$ is the configuration entropy of the $q$-field.
$\int|\mathbf{J}|^2/(2\kappa)\,d^3x$ is the kinetic energy of information currents.
Damping $\gamma_0(1-q)$ dissipates kinetic energy at a rate proportional to local occupation, driving the system toward the entropy-maximizing static configuration.

\subsection{Time-Dependent Stability Analysis}
\label{sec:stability}

The self-organization decomposition~\eqref{eq:decomp} describes the static balance.
For time-dependent configurations, the full telegraph equation governs the competition between diffusion and collapse.
Consider a localized perturbation with characteristic scale $L$ and amplitude $\delta q$.
Three timescales compete:

\begin{enumerate}
    \item \textbf{Wave timescale:} $\tau_{\rm wave} \sim L/c$ --- the light-crossing time.
    \item \textbf{Damping timescale:} $\tau_{\rm damp} \sim 1/[\gamma_0(1-q_0)]$ --- the time for damping to dissipate coherent oscillations.
    \item \textbf{Diffusion timescale:} $\tau_{\rm diff} \sim L^2/D_{\rm eff} = L^2\gamma_0(1-q_0)/c^2$ --- the time for diffusion to smooth gradients.
\end{enumerate}

The dimensionless ratio controlling the dynamics is:

\begin{equation}
R \equiv \frac{\tau_{\rm damp}}{\tau_{\rm wave}} = \frac{c}{\gamma_0(1-q_0)L}.
\label{eq:stability_R}
\end{equation}

\begin{itemize}
    \item $R\gg1$ (small $L$, high $q_0$): wave propagation dominates. Perturbations oscillate and disperse.
    \item $R\ll1$ (large $L$, low $q_0$): diffusion dominates. Perturbations smooth out on timescale $\tau_{\rm diff}$.
    \item The self-organization term $|\nabla q|^2/q^2$ becomes competitive with diffusion when $|\nabla q|/q \gtrsim 1/L$, i.e., when the logarithmic gradient exceeds the inverse scale. This occurs near mass concentrations where $q\ll1$.
\end{itemize}

A perturbation is unconditionally stable against gravitational collapse when $R\gg1$---the wave dynamics disperse the perturbation before self-organization can amplify it.
Conversely, when $R\ll1$ \emph{and} $q\ll1$, the self-organization term dominates both wave dispersion and diffusion, driving runaway collapse.
This provides a quantitative criterion for the onset of gravitational instability in the $q$-field framework.

\subsection{Variational Interpretation}

Equation~\eqref{eq:laplace_ln} can be understood variationally.
Define the information-theoretic action:

\begin{equation}
\mathcal{S}[q] = \int d^3x \, \frac{|\nabla q|^2}{q^2}.
\label{eq:action}
\end{equation}

The Euler--Lagrange equation $\delta\mathcal{S}/\delta q=0$ yields $\nabla^2(\ln q)=0$.
The exponential solution~\eqref{eq:q_solution} extremizes the integrated squared logarithmic gradient.
Physically: nature settles into the configuration that minimizes information gradients for a given total information charge (mass).
The $1/r$ potential is not a force law but a \textbf{minimum free-energy configuration of the $q$-field}, as proven by the H-theorem of Sec.~\ref{sec:lyapunov}.

\subsection{Quantitative vs.\ Qualitative Predictions}
\label{sec:qual_vs_quant}

We must distinguish predictions that depend on the specific choice $\phi(q)=-\ln q$ from those protected by universality.

\paragraph{Qualitative predictions (universal across $\Phi$):}
\begin{itemize}
    \item Damping vanishes in vacuum: $\gamma(q)\to0$ as $q\to1$.
    \item Lorentz invariance is dynamically selected by vacuum.
    \item The static equation is $\nabla^2\phi(q)=0$; the $1/r$ potential emerges for any $\phi\in\Phi$.
    \item Diffusion and self-organization are two regimes of a single equation.
    \item The wave--diffusion crossover exists at $k=\gamma/(2c)$.
    \item A soft boundary replaces the hard event horizon.
    \item The H-theorem holds (with a $\phi$-dependent Lyapunov functional).
\end{itemize}

\paragraph{Quantitative predictions (specific to $\phi(q)=-\ln q$):}
\begin{itemize}
    \item The exact static solution: $q(r)=e^{-GM/rc^2}$.
    \item The Newtonian limit: $\Phi(r)=-GM/r + O(r^{-2})$.
    \item The specific form of the matter--$q$-field coupling (source term).
    \item The exact value of $\gamma_0=c/\ell_P$ (depends on the naturalness argument, not $\phi$).
\end{itemize}

Observational tests of quantitative predictions can distinguish $\phi(q)=-\ln q$ from other members of $\Phi$, while qualitative predictions test the entire universality class.

%%%%%%%%%%%%%%%%%%%%%%%%%%%%%%%%%%%%%%%%%%%%%%%%%%%%%%%%%%%%%%%%%%%%%%
\section{Testable Predictions}
\label{sec:predictions}
%%%%%%%%%%%%%%%%%%%%%%%%%%%%%%%%%%%%%%%%%%%%%%%%%%%%%%%%%%%%%%%%%%%%%%

\subsection{Wave--Diffusion Crossover}

The most distinctive prediction is a frequency-dependent crossover from wave propagation to diffusion in gravitational dynamics.
From the dispersion relation of Eq.~\eqref{eq:telegraph_gammaq} for a plane-wave perturbation $q\propto e^{i(kx-\omega t)}$ on a background $q_0$:

\begin{equation}
\omega^2 + i\gamma_0(1-q_0)\omega = c^2 k^2/q_0,
\label{eq:dispersion}
\end{equation}

the crossover wavenumber is:

\begin{equation}
\boxed{k_{\rm cross} = \frac{\gamma_0(1-q_0)\sqrt{q_0}}{2c}}.
\label{eq:kcross}
\end{equation}

For $k\gg k_{\rm cross}$ (short wavelengths): the $\partial_t^2 q$ term dominates, and $q$ propagates as a damped wave at speed $c$---conventional gravitational wave behavior.
For $k\ll k_{\rm cross}$ (long wavelengths): the $\gamma\partial_t q$ term dominates, and $q$ evolves diffusively---classical gravitational relaxation.

Observable consequences:
\begin{itemize}
    \item \textbf{Frequency-dependent dispersion.} The phase velocity $v_{\rm phase}=c/\sqrt{1+i\gamma/\omega}$ deviates from $c$ at frequencies near and below $\omega_{\rm cross}=\gamma/2$. For a broadband gravitational wave signal (e.g., a binary inspiral), this produces frequency-dependent arrival time delays.

    \item \textbf{Suppression of low-frequency GWs.} Modes with $\omega<\gamma/2$ are overdamped---they decay exponentially without oscillation.
    This suppresses gravitational wave emission below a cutoff frequency, affecting the low-frequency end of the stochastic gravitational wave background.

    \item \textbf{Ringdown deviations.} The quasi-normal mode spectrum of a post-merger compact object is modified: modes above $k_{\rm cross}$ oscillate; modes below decay monotonically.
\end{itemize}

With $\gamma_0=c/\ell_P\sim10^{43}$ s$^{-1}$, the crossover frequency is Planckian and irrelevant for observations.
However, the \emph{effective} damping $\gamma=\gamma_0(1-q)$ is far smaller in realistic environments.
Near a solar-mass object at $r\sim R_\odot$, $\gamma\sim\gamma_0\times2\times10^{-6}\sim4\times10^{37}$ s$^{-1}$, still astronomically large.
For the crossover to enter the observable band (e.g., LISA frequencies $\sim10^{-3}$ Hz), we would need $(1-q)\sim\gamma/\gamma_0\sim10^{-46}$, which occurs only in extremely deep vacuum far from all mass.

This means the \textbf{wave regime} is the generic prediction for all observable gravitational wave phenomena.
The diffusion regime is relevant only in the extreme strong-field regime ($q\ll1$), where it may affect black hole interiors, the very early universe, or the end state of gravitational collapse.

\textbf{Honest statement:} With $\gamma_0$ fixed by the Planck scale, the wave--diffusion crossover is not observable with current or planned detectors.
Its significance is theoretical: it provides a conceptual unification of wave (quantum coherent) and diffusive (classical decoherent) gravitational dynamics within a single equation. Whether this unification is physically meaningful or merely formal is, in the end, a question for experiments we do not yet know how to do.

\subsection{Soft Boundary: No Hard Event Horizon}

In Eq.~\eqref{eq:q_solution}, $q(r)>0$ for all finite $r$:

\begin{equation}
q(r) > 0 \quad \text{for all} \quad r > 0.
\label{eq:soft}
\end{equation}

There is \textbf{no event horizon} in the strict sense.
The free capacity vanishes only asymptotically as $r\to0$.
Matter falling toward $r=0$ encounters an ever-decreasing $q$, approaching but never reaching the capacity bound.
A \textbf{soft boundary}. Gradual saturation of information capacity without sharp causal disconnection.

The Event Horizon Telescope observations of M87*~\cite{EventHorizonTelescope:2019dse} and Sgr A*~\cite{EventHorizonTelescope:2022xnr} reveal dark central shadows consistent with a photon sphere at $\sim2.6\,GM/c^2$ (for M87*).
In the $q$-field picture, the shadow is produced not by an event horizon but by the steep gradient of $q$ near the origin, which traps photons gravitationally.
The profile may differ from the Kerr metric prediction at the percent level.

Distinguishing the soft boundary from the hard Kerr horizon requires:
\begin{itemize}
    \item Higher-resolution EHT observations (next-generation EHT, space VLBI).
    \item Time-domain observations of orbiting hot spots near the ISCO.
    \item Gravitational wave ringdown spectra from binary black hole mergers, whose quasi-normal mode frequencies may differ between a hard horizon and a soft boundary.
\end{itemize}

\subsection{Gravitational Wave Speed and the Milky Way Halo}
\label{sec:gw170817_halo}

In the vacuum limit ($q\to1$, $\gamma\to0$), gravitational waves propagate at exactly $c$.
This is consistent with GW170817/GRB 170817A~\cite{Abbott:2017apx,GW170817grb}, which constrain $-3\times10^{-15}\leq(v_{\rm GW}-c)/c\leq7\times10^{-16}$.

\textbf{Galactic halo propagation: an open theoretical question.}
GW170817 propagated through the Milky Way halo before detection.
In the $q$-field framework, matter-filled regions have $q<1$ and therefore $\gamma>0$.
With $1-q_0 \sim 5\times10^{-7}$ at the Solar position ($r\approx8$\,kpc, $M_{\rm enc}\approx10^{11}M_\odot$), the effective damping is $\gamma_{\rm halo} \approx 9\times10^{36}$\,s$^{-1}$, so $\gamma_{\rm halo}/\omega \sim 10^{34}$ at LIGO frequencies.
In this strongly overdamped regime, freely propagating scalar $q$-field modes ($k < k_{\rm cross}$) have purely imaginary eigenfrequencies---they decay exponentially in time without oscillation.
No coherent scalar signal propagates through the halo.

\textbf{Conjectured resolution.}
If the gravitational wave signal detected by LIGO corresponds to \emph{tensor} perturbations $h_{\mu\nu}$ of an effective metric $g_{\mu\nu}[q]$ that emerges from the static $q$-field background, and if tensor modes decouple from the scalar sector's damping at linear order, then $\Box h_{\mu\nu}=0$ in vacuum and tensor waves propagate at $c$ regardless of the scalar sector's behavior.
This decoupling is a conjecture requiring future proof: the paper has not derived the effective metric or the tensor perturbation equations from the $q$-field.
Until such a derivation exists, the consistency of the $q$-field framework with GW170817 must be regarded as an open theoretical question, not a resolved check.

\subsection{Separation of Qualitative and Quantitative Predictions}

Following the classification of Sec.~\ref{sec:qual_vs_quant}, we separate:

\begin{center}
\begin{tabular}{p{0.45\columnwidth}p{0.45\columnwidth}}
\toprule
\textbf{Qualitative (universal across $\Phi$)} & \textbf{Quantitative ($\phi=-\ln q$ specific)} \\
\midrule
Wave--diffusion crossover exists & Crossover wavenumber $k_{\rm cross}=\gamma/(2c)$ \\
Soft boundary replaces horizon & Exact $q(r)=e^{-GM/rc^2}$ profile \\
Lorentz selected by vacuum & $\gamma_0=c/\ell_P$ (naturalness) \\
$1/r$ potential emerges (any $\phi$) & Newtonian limit $\Phi=-GM/r$ \\
Damping vanishes at $q\to1$ & $v_{\rm GW}/c = 1-(1-q_0)^2$ \\
\bottomrule
\end{tabular}
\end{center}

%%%%%%%%%%%%%%%%%%%%%%%%%%%%%%%%%%%%%%%%%%%%%%%%%%%%%%%%%%%%%%%%%%%%%%
\section{Discussion}
\label{sec:discussion}
%%%%%%%%%%%%%%%%%%%%%%%%%%%%%%%%%%%%%%%%%%%%%%%%%%%%%%%%%%%%%%%%%%%%%%

\subsection{Summary of the Framework}

Two axioms. One structural choice. The results, in compressed form:

\begin{enumerate}
    \item The unified current equation~\eqref{eq:current_fundamental} with logarithmic gradient and $\bar{q}$-dependent damping yields the telegraph equation~\eqref{eq:telegraph_unified}: $\partial_t^2 q + \gamma_0(1-q)\partial_t q = c^2\nabla^2(\ln q)$. This one equation governs all regimes---no patch, no separate linear and nonlinear equations.
    \item Lorentz invariance is dynamically selected by vacuum ($q\to1$); preferred-frame effects are confined to matter-filled regions.
    \item The static limit follows directly from $\partial_t J_{ij}=0$: the factor $\gamma_0(1-\bar{q})$ cancels, yielding $\nabla^2(\ln q)=0$, whose spherically symmetric solution $q(r)=e^{-GM/rc^2}$ recovers Newtonian gravity.
    \item The same equation unifies decoherence and gravitational collapse via $\nabla^2 q/q = |\nabla q|^2/q^2$.
    \item An H-theorem is proved (Sec.~\ref{sec:lyapunov}): $dE/dt\leq0$ for the Lyapunov functional $E=\int[\Phi(q)+|\mathbf{J}|^2/(2\kappa)]d^3x$, establishing thermodynamic consistency.
    \item The theory predicts a wave--diffusion crossover and a soft boundary at black holes.
\end{enumerate}

\subsection{The Holographic Tension}
\label{sec:holographic}

Now the hard part.
The capacity bound of 1 bit per cell is a \textbf{local} bound applied to each cell independently.
For a region of size $R$, the total information capacity is proportional to the number of cells in that region, scaling as the volume:

\begin{equation}
S_{\rm max} \sim \left(\frac{R}{a}\right)^3 \propto R^3 \qquad \text{(Volume scaling)}.
\label{eq:volume_scaling}
\end{equation}

This contrasts with the holographic principle~\cite{tHooft:1993,Susskind:1994vu}, which posits that the maximum entropy of a region scales with its surface area:

\begin{equation}
S_{\rm holographic} \sim \frac{A}{4\ell_P^2} \propto R^2 \qquad \text{(Area scaling)}.
\label{eq:area_scaling}
\end{equation}

Volume versus area. This is a qualitative difference in the scaling of maximum information content.

\textbf{Why this matters.}
The holographic principle is supported by two independent lines of evidence:
(1) black hole thermodynamics, where $S_{\rm BH}=A/4\ell_P^2$ is proportional to horizon area, and
(2) the AdS/CFT correspondence, where a gravitational theory in $d+1$ dimensions is equivalent to a nongravitational theory in $d$ dimensions.
If the $q$-field theory is to be taken seriously as a description of gravity, this tension must be resolved.

\textbf{A conjectured resolution.}
A natural pathway out of this tension is that the 1-bit-per-cell bound is an \emph{ultraviolet} capacity limit, while the \emph{infrared} entropy scaling is determined by entanglement structure---not by how many cells exist in a volume, but by how many are pairwise entangled across its boundary.
If each cell within one Planck length of the boundary is maximally entangled with a cell outside, the entanglement entropy scales as the boundary area: $S_{\rm ent}\sim(R/a)^2\ln 2\propto R^2$.
This resolution is a hypothesis requiring formal proof within the $q$-field dynamics; establishing it rigorously is a priority for future work.

\subsection{Fair Comparison with HTMAU Using Tiered Methodology}
\label{sec:htmau_comparison}

Bassani and Magueijo (2025)~\cite{Bassani:2025htma} proposed ``How to Make a Universe'' (HTMAU).
In their framework, diffeomorphism invariance emerges from a Markov chain of stochastic choices evolving toward an absorbing state.
Both HTMAU and the present $q$-field theory share a common conviction: spacetime symmetries are emergent rather than fundamental.
The inputs, mechanisms, and outputs differ substantially, however.
To make the comparison systematic, we apply the \textbf{same tiered methodology} (Sec.~\ref{sec:tiered}) to both theories.

\subsubsection{HTMAU's Assumptions in Tiered Classification}

Enumerating HTMAU's specific assumptions as presented in Ref.~\cite{Bassani:2025htma}:

\begin{enumerate}
    \item \textbf{Preferred foliation.} Spacetime has a preferred time foliation---a global time coordinate that singles out a rest frame. This is an irreducible structural input (Tier~0 equivalent).
    \item \textbf{Global/local variable split.} The universe's state is partitioned into ``global'' variables (constants of nature) and ``local'' variables (field configurations on the foliation). This is a structural choice (Tier~2 equivalent).
    \item \textbf{Matter-gravity separation.} The total action separates into a gravitational sector and a matter sector, with the matter sector treated as a perturbation. This is a modeling choice (Tier~2 equivalent).
    \item \textbf{Local diffeomorphism invariance is imposed before the absorbing state.} The Markov chain is constructed to \emph{preserve} local spatial diffeomorphisms at each step, with time reparameterization invariance emerging only in the absorbing state. This means diffeomorphism invariance is partially \emph{assumed}, not fully derived (Tier~1 equivalent---an empirical input about the symmetry structure).
    \item \textbf{Constants as time-conjugate variables.} The ``global'' variables are canonically conjugate to time, with Poisson bracket structure imposed. This is a structural choice (Tier~2 equivalent).
    \item \textbf{Beta potential functions.} Two $\beta$ potential functions ($\beta_{\rm local}$ and $\beta_{\rm global}$) control the mutation rule of the Markov chain. Their specific forms are free functions (Tier~2 structural choices with continuous parameters).
    \item \textbf{Markov chain mutation rule.} The specific mutation rule $X\to X+\{X,H\}_{\rm local}\delta t + \{\beta_{\rm global},X\}\delta W$ is a structural ansatz (Tier~2 equivalent).
\end{enumerate}

\subsubsection{Tiered Comparison}

Table~\ref{tab:compare_htmau} compares the two frameworks using the same tiered methodology.

\begin{table}[h]
\centering
\caption{Tiered comparison of emergent-spacetime frameworks. Both theories are classified using the same four-tier methodology (Sec.~\ref{sec:tiered}).}
\label{tab:compare_htmau}
\begin{tabular}{p{0.22\columnwidth}p{0.35\columnwidth}p{0.35\columnwidth}}
\toprule
& \textbf{HTMAU (Bassani \& Magueijo 2025)} & \textbf{This work ($q$-field)} \\
\midrule
\textbf{Tier 0} & Causality (via Markov ordering) & A1: Causality \\
(Physical & & A2: 1-bit capacity bound \\
Axioms) & Preferred foliation & \\
& (2 axioms) & (2 axioms) \\
\midrule
\textbf{Tier 1} & $c$, $\hbar$, $G$ & $c$, $\hbar$, $G$, $d=3$ \\
(Empirical & Local diffeo invariance & \\
Inputs) & (imposed pre-absorbing state) & \\
& (4 inputs) & (4 inputs) \\
\midrule
\textbf{Tier 2} & Global/local variable split & Regular 3D lattice \\
(Modeling & Matter-gravity separation & Nearest-neighbor edges \\
Choices) & Constants as time-conjugates & $\phi(q)=-\ln q\in\Phi$ \\
& 2 $\beta$ potential functions & \\
& Markov mutation rule & \\
& (5 structural choices) & (3 structural choices) \\
\midrule
\textbf{Tier 3} & Mutation rate parameters & $\gamma_0=c/\ell_P$ \\
(Derived) & $\beta$ function parameters & $D_0=c^2/\gamma_0=c\ell_P$ \\
& (2+ free parameters) & (0 additional parameters) \\
\midrule
\textbf{Emergent} & Diffeomorphism invariance & Lorentz invariance \\
\textbf{symmetry} & (full GR symmetry) & (SR symmetry) \\
\midrule
\textbf{Emergence} & Stochastic: Markov chain & Deterministic: \\
\textbf{mechanism} & $\to$ absorbing state & $q\to1 \Rightarrow \gamma\to0$ \\
\midrule
\textbf{Static} & FRW matter production rate & $\nabla^2(\ln 1/q)=0$ \\
\textbf{equation} & (cosmological) & $q(r)=e^{-GM/rc^2}$ \\
\midrule
\textbf{Testable} & Near-diffeo invariant & Wave--diffusion crossover, \\
\textbf{predictions} & cosmology & soft boundary, $v_{\rm GW}$ correction \\
\bottomrule
\end{tabular}
\end{table}

\subsubsection{Interpretation}

Both frameworks share a philosophy of emergent spacetime.
HTMAU achieves the stronger result---diffeomorphism invariance, the full symmetry of GR---but requires substantially more input: a preferred foliation (an additional Tier~0 axiom), local diffeomorphism invariance imposed pre-emergence (a Tier~1 input), and five Tier~2 structural choices including two free $\beta$ functions.
The $q$-field framework achieves the weaker result---Lorentz invariance, sufficient for special relativity and the Newtonian limit---but with fewer assumptions at every tier.

Is diffeomorphism invariance necessary at the fundamental level, or is it sufficient to show that Lorentz invariance emerges and the Newtonian limit is recovered?
If the latter, the economy of assumptions favors the present approach.
If the former, HTMAU's richer structure is necessary and the present theory must be extended---a direction discussed in the open directions below.

We emphasize: this comparison is methodological, not adversarial.
Both approaches contribute to understanding spacetime as emergent.
The tiered methodology provides a systematic framework for comparing economy and explanatory power across different approaches.

\subsection{Honest Assessment of Limitations}

\begin{enumerate}
    \item \textbf{$\gamma(q)=\gamma_0(1-q)$ is a modeling ansatz.} The linear form is the simplest model consistent with A2 but is not rigorously derived from it. Nonlinear corrections $\gamma\propto(1-q)^\alpha$ with $\alpha\neq1$ are possible. The essential prediction---$\gamma\to0$ as $q\to1$---is robust regardless of the functional form.

    \item \textbf{The choice $\phi(q)=-\ln q$ is protected by universality but not unique.} Any $\phi\in\Phi$ gives the same qualitative physics. The $-\ln q$ form is the simplest and most natural, but the universality argument must be made rigorous through renormalization-group analysis.

    \item \textbf{We recover only the Newtonian limit of GR.} The full Einstein equations, with their nonlinear self-interaction, frame-dragging, and cosmological constant, are not derived. Extending the $q$-field framework to full GR requires deriving the right-hand side of $\nabla^2(\ln 1/q)=(\text{source term})$ for general stress-energy configurations, and showing that the result matches the Einstein tensor in the appropriate limit.

    \item \textbf{Matter coupling is not fully derived.} We assumed spherical symmetry and identified an integration constant with $GM/c^2$. A complete theory must derive how $T_{\mu\nu}$ sources $q$ in general configurations. The matter--$q$-field coupling proposed in Sec.~\ref{sec:static} is a first step but requires rigorous derivation from the microscopic dynamics.

    \item \textbf{Spatial dimensionality is an input.} $d=3$ is taken as an empirical fact, not a prediction. The $1/r$ potential follows from $d=3$; in $d$ dimensions the static solution would be $q\propto\exp(-\text{const}/r^{d-2})$, yielding a $1/r^{d-2}$ force law.

    \item \textbf{The holographic bound is not satisfied.} As discussed in Sec.~\ref{sec:holographic}, the 1-bit-per-cell bound implies volume scaling $S\sim R^3$, contradicting the holographic area scaling $S\sim R^2$. This is a serious limitation that requires resolution through nonlocal entanglement, holographic cell arrangement, or dynamical saturation constraints.

    \item \textbf{Lorentz selection is weaker than diffeomorphism emergence.} Our theory always contained a Lorentz-invariant sector; the result is that vacuum selects it. HTMAU's diffeomorphism emergence is a stronger result. Whether this matters depends on whether full diffeomorphism invariance is needed at the fundamental level or only in the classical gravitational limit.

    \item \textbf{Tensor vs.\ scalar sector distinction.} The resolution of the GW170817 halo tension (Sec.~\ref{sec:gw170817_halo}) relies on a distinction between scalar $q$-field perturbations (subject to damping) and tensor metric perturbations (propagating at $c$). A complete theory must derive the coupling between the scalar and tensor sectors from first principles, showing that the decoupling at linear order is robust.
\end{enumerate}

\subsection{Open Directions}

\begin{enumerate}
    \item \textbf{Full GR recovery.} The most important open problem is extending the $q$-field framework to recover the full Einstein equations. The nonlinear structure of $\nabla^2(\ln 1/q)=0$ hints at a deeper connection to the Ricci tensor, but the mapping is not yet known. A Poisson-type equation for $\ln q$ sourced by $\propto(1-q)/a^3$ may relate to the $G_{00}$ component of the Einstein tensor in the Newtonian gauge.

    \item \textbf{Holographic resolution.} Resolving the volume-vs.-area scaling tension (Sec.~\ref{sec:holographic}) is essential. This may involve replacing the regular 3D lattice with a holographic cell arrangement, or deriving an effective area-law bound from the dynamics of entangled cells.

    \item \textbf{Cosmology.} Global expansion corresponds to $q\to1$ (increasing free capacity), which is natural if the universe's information content grows more slowly than its capacity. This may connect to dark energy. The $q$-field should have cosmological solutions that can be compared with $\Lambda$CDM.

    \item \textbf{Entanglement and ER = EPR.} Zilly's QM foundation involves entanglement. In the $q$-field picture, two entangled cells share information, modifying their local $q$. The connection between entanglement structure and spacetime geometry~\cite{Maldacena:2013je} may be derivable in this framework.

    \item \textbf{Quantum gravity regime.} When $q\to0$, the capacity bound is nearly saturated. Zilly's finite-capacity QM predicts strong quantum effects. The $q$-field dynamics in this regime may describe quantum gravitational phenomena without a full theory of quantum gravity.

    \item \textbf{Numerical simulations.} Simulating the $q$-field dynamics in binary merger scenarios and comparing waveforms with GR predictions. The soft boundary and nonlinear damping should produce distinguishable ringdown signals.

    \item \textbf{Experimental constraints on $\gamma_0$ and $\phi(q)$.} While we fix $\gamma_0=c/\ell_P$ from naturalness and $\phi(q)=-\ln q$ from simplicity, both choices are falsifiable. If observations constrain $\gamma_0$ to be smaller (e.g., from the absence of anomalous damping in binary pulsar timing), the identification with the Planck scale would need revision. Similarly, deviations from the $e^{-GM/rc^2}$ profile would constrain $\phi(q)$.

    \item \textbf{Rigorous universality proof.} The claim that any $\phi\in\Phi$ gives the same qualitative physics must be made rigorous. This requires a renormalization-group analysis showing that $\phi(q)=-\ln q$ is the infrared fixed point of the universality class, or a constructive proof that macroscopic observables (potential, force law, stability) are independent of the microscopic $\phi$.
\end{enumerate}

%%%%%%%%%%%%%%%%%%%%%%%%%%%%%%%%%%%%%%%%%%%%%%%%%%%%%%%%%%%%%%%%%%%%%%
\section*{Acknowledgments}
%%%%%%%%%%%%%%%%%%%%%%%%%%%%%%%%%%%%%%%%%%%%%%%%%%%%%%%%%%%%%%%%%%%%%%

We thank J.~Zilly for providing the foundational proof that finite information capacity plus self-referential consistency yields quantum mechanics~\cite{Zilly:2026}, without which this work would lack its axiomatic grounding.
We thank participants of the Emergent Gravity workshop for discussions.
We thank Bassani and Magueijo for making their HTMAU manuscript available, which provided a valuable benchmark against which to measure the economy of our own assumptions and motivated the tiered classification methodology.

%%%%%%%%%%%%%%%%%%%%%%%%%%%%%%%%%%%%%%%%%%%%%%%%%%%%%%%%%%%%%%%%%%%%%%
\bibliographystyle{apsrev4-2}
\begin{thebibliography}{99}

\bibitem{Zilly:2026}
J.~Zilly, ``Existence as Distinguishability: Quantum Mechanics from Finite Graded Equality,''
arXiv:2603.11900 (2026).
[arXiv preprint; independent verification pending. Not re-derived in this work.]

\bibitem{Neukart:2024qmm}
F.~Neukart, R.~Brasher, and E.~Marx,
``The Quantum Memory Matrix: A Unified Framework for the Black Hole Information Paradox,''
Entropy \textbf{26}, 1039 (2024).
[Independent validation of the Planck-cell paradigm.]

\bibitem{Green:1987sp}
M.~B.~Green, J.~H.~Schwarz, and E.~Witten,
\textit{Superstring Theory},
Cambridge University Press (1987).

\bibitem{Polchinski:1998rq}
J.~Polchinski,
\textit{String Theory},
Cambridge University Press (1998).

\bibitem{Rovelli:2004tv}
C.~Rovelli,
\textit{Quantum Gravity},
Cambridge University Press (2004).

\bibitem{Ashtekar:2004eh}
A.~Ashtekar and J.~Lewandowski,
``Background independent quantum gravity: A status report,''
Class.\ Quant.\ Grav.\ \textbf{21}, R53 (2004).

\bibitem{Ambjorn:2001cv}
J.~Ambjorn, J.~Jurkiewicz, and R.~Loll,
``Dynamically triangulating Lorentzian quantum gravity,''
Nucl.\ Phys.\ B \textbf{610}, 347 (2001).

\bibitem{Weinberg:1979}
S.~Weinberg,
``Ultraviolet divergences in quantum theories of gravitation,''
in \textit{General Relativity: An Einstein Centenary Survey},
Cambridge University Press (1979).

\bibitem{Reuter:2012}
M.~Reuter and F.~Saueressig,
``Quantum Einstein gravity,''
New J.\ Phys.\ \textbf{14}, 055022 (2012).

\bibitem{Verlinde:2010hp}
E.~P.~Verlinde,
``On the origin of gravity and the laws of Newton,''
JHEP \textbf{04}, 029 (2011).

\bibitem{Jacobson:1995ab}
T.~Jacobson,
``Thermodynamics of spacetime: The Einstein equation of state,''
Phys.\ Rev.\ Lett.\ \textbf{75}, 1260 (1995).

\bibitem{Padmanabhan:2010}
T.~Padmanabhan,
``Thermodynamical aspects of gravity: New insights,''
Rep.\ Prog.\ Phys.\ \textbf{73}, 046901 (2010).

\bibitem{Bassani:2025htma}
P.~M.~Bassani and J.~Magueijo,
``How to make a universe,''
Phys.\ Rev.\ D \textbf{111}, 103529 (2025).

\bibitem{Cattaneo:1948}
C.~Cattaneo,
``Sulla conduzione del calore,''
Atti Sem.\ Mat.\ Fis.\ Univ.\ Modena \textbf{3}, 83 (1948).

\bibitem{Jou:2010}
D.~Jou, J.~Casas-V{\'a}zquez, and G.~Lebon,
\textit{Extended Irreversible Thermodynamics},
Springer (2010).

\bibitem{Abbott:2017apx}
B.~P.~Abbott \textit{et al.} (LIGO Scientific and Virgo Collaborations),
``GW170817: Observation of gravitational waves from a binary neutron star inspiral,''
Phys.\ Rev.\ Lett.\ \textbf{119}, 161101 (2017).

\bibitem{GW170817grb}
B.~P.~Abbott \textit{et al.},
``Gravitational waves and gamma-rays from a binary neutron star merger: GW170817 and GRB 170817A,''
Astrophys.\ J.\ Lett.\ \textbf{848}, L13 (2017).

\bibitem{EventHorizonTelescope:2019dse}
Event Horizon Telescope Collaboration,
``First M87 Event Horizon Telescope results. I. The shadow of the supermassive black hole,''
Astrophys.\ J.\ Lett.\ \textbf{875}, L1 (2019).

\bibitem{EventHorizonTelescope:2022xnr}
Event Horizon Telescope Collaboration,
``First Sagittarius A* Event Horizon Telescope results. I. The shadow of the supermassive black hole in the center of the Milky Way,''
Astrophys.\ J.\ Lett.\ \textbf{930}, L12 (2022).

\bibitem{Bassi:2013mta}
A.~Bassi, K.~Lochan, S.~Satin, T.~P.~Singh, and H.~Ulbricht,
``Models of wave-function collapse, underlying theories, and experimental tests,''
Rev.\ Mod.\ Phys.\ \textbf{85}, 471 (2013).

\bibitem{Maldacena:2013je}
J.~Maldacena and L.~Susskind,
``Cool horizons for entangled black holes,''
Fortsch.\ Phys.\ \textbf{61}, 781 (2013).

\bibitem{Will:2014}
C.~M.~Will,
``The confrontation between general relativity and experiment,''
Living Rev.\ Relativ.\ \textbf{17}, 4 (2014).

\bibitem{Kostelecky:2008}
V.~A.~Kosteleck{\'y} and N.~Russell,
``Data tables for Lorentz and CPT violation,''
Rev.\ Mod.\ Phys.\ \textbf{83}, 11 (2011).

\bibitem{Mewes:2013}
M.~Mewes,
``Tests of Lorentz invariance,''
in \textit{Proceedings of the 6th Meeting on CPT and Lorentz Symmetry},
World Scientific (2013).

\bibitem{tHooft:1993}
G.~'t Hooft,
``Dimensional reduction in quantum gravity,''
arXiv:gr-qc/9310026 (1993).

\bibitem{Susskind:1994vu}
L.~Susskind,
``The world as a hologram,''
J.\ Math.\ Phys.\ \textbf{36}, 6377 (1995).

\end{thebibliography}

\end{document}
